\documentclass[10pt,a4paper]{amsart}
\usepackage[margin=19mm]{geometry}
\usepackage{amsmath,amssymb,mathtools}
\usepackage[scr=rsfs,cal=euler]{mathalfa}
\usepackage{xfrac}
\usepackage[unicode,hidelinks,bookmarks=false]{hyperref}
\newcommand{\ie}{i.e.\ }
\newcommand{\cf}{cf.\ }
\newcommand{\resp}{respectively\ }
\newcommand{\Subsection}{\subsection}
\providecommand{\nobreakdash}{\nobreak\hbox{-}}
\setlength{\emergencystretch}{2em}
\multlinegap0pt
\usepackage{bm}
\usepackage{bbm}
\usepackage{dsfont}
\usepackage{xspace}
\usepackage{cancel}
\usepackage{mathtools}
\usepackage{amsthm}
\makeatletter
\@ifundefined{theorem}{\newtheorem{theorem}{Theorem}}{}
\makeatother
\makeatletter
\expandafter\ifx\csname proof\endcsname\relax
\newenvironment{proof}{\noindent{\bf Proof:}}{\qed}
\fi
\makeatother
\title[Permutations with a fixed number of occurrences of a monoton]{Permutations with a fixed number of occurrences of a monotone pattern}
\author{Michael Waite}
\date{Source: arXiv:2506.05712v2 (CC BY 4.0)}
\begin{document}
\maketitle

\noindent\emph{Source and licence.} Permutations with a fixed number of occurrences of a monotone pattern. Version arXiv:2506.05712v2 (CC BY 4.0), distributed under the Creative Commons Attribution 4.0 International licence, as recorded in the arXiv licence metadata for this exact version and shown on the abstract page https://arxiv.org/abs/2506.05712v2. Full rights notice: licenses/arxiv-2506.05712-CC-BY-4.0.txt.\par\smallskip

\noindent\textit{Editorial scope.} Counted unit: the Theorem whose source label is \texttt{None} in the article (source0.tex lines 752--757, source label \texttt{None}), delivered together with the article's own complete original proof of it (source0.tex lines 758--776, one \texttt{proof} environment of 19 lines). No mathematics is written, repaired or re-derived: statement and proof are the article's own text line for line. Required context retained uncounted: none. Every definition, display and label of those lines is retained unchanged. Omitted from this package and not redistributed: 1--751, 777--880. Nothing inside a retained excerpt is omitted or altered: every retained body is delivered exactly as the article's lines are written, so no omission receipt of any kind accompanies this package. Citations: the retained excerpts carry no citation command at all, so no bibliography entry of the article is transcribed and no reference is redistributed. Reference adaptation: none was needed and none was made; every reference inside the delivered unit resolves inside it, so no unresolved reference or citation is produced. Printed slips kept literal: no printed word, symbol, hypothesis or number of the counted statement or of its proof was altered. Layout and class: the article's own class is \texttt{amsart} with packages including graphicx, latexsym, amsfonts, amsmath, amssymb, url, tikz, float, color; the wrapper is the lane's pinned \texttt{amsart} editorial wrapper at 10pt on A4, which restates the theorem-like environments the retained text needs and adds the article's own macro definitions that the retained text needs (none), with extra wrapper packages (bm, bbm, dsfont, xspace, cancel, mathtools). The delivered numbering is the unit's own. Assets: no retained span contains a figure environment, a graphics-inclusion command, a PDF-page inclusion, a table or a TikZ picture; no image file is copied into this package, the package declares no asset dependency and no image file is redistributed with it. Licence: version arXiv:2506.05712v2 of this article is distributed under the Creative Commons Attribution 4.0 International licence, as recorded in the arXiv licence metadata for this exact version and shown on the abstract page https://arxiv.org/abs/2506.05712v2. A full rights notice is delivered at licenses/arxiv-2506.05712-CC-BY-4.0.txt. Characters: every retained excerpt is pure ASCII, so the delivered engine, XeLaTeX with its default Latin Modern text font, reports no missing character.
\begin{theorem}
    Let $p_0$ be a permutation of any size (including $0$). Then there is a constant $K_r$ so that 
    \[
    |S_{n,r}(321\ominus p_0)| \leq K_r |S_{n}(321\ominus p_0)|.
    \]
\end{theorem}

\begin{proof}
    We know that for each possible $q$, we have
    \[
    |S^q_{n,r}(321\ominus p_0)| \leq |S_{n + m_{\overline{q}}}(231\ominus p_0)|.
    \]
    Applying Lemma $3.4$, this gives us 
    \[
    |S^q_{n,r}(321\ominus p_0)| \leq |S_{n + m_{\overline{q}}}(321\ominus p_0)|.
    \]
    Now $|S_{n + m_{\overline{q}}}(321\ominus p_0)| \leq K^{(1)}_r |S_{n}(321\ominus p_0)| $ for some constant $K^{(1)}_r$. So
    \[
    |S^q_{n,r}(321\ominus p_0)| \leq K^{(1)}_r |S_{n}(321\ominus p_0)|.
    \]
    Since $q$ is a permutation that has $r$ copies of $321\ominus p_0$ where every entry of $q$ participates in a $321\ominus p_0$ pattern, then $|q| \leq r(3 + |p_0|)$. So there are at most $1! + 2! + ... + (r(3 + |p_0|))!$ choices for $q$. Set $K^{(2)}_r = 1! + 2! + ... + (r(3 + |p_0|))!$. Now we get
    \[
    |S_{n,r}(321\ominus p_0)| \leq K^{(1)}_r K^{(2)}_r |S_{n}(321\ominus p_0)|.
    \]
    Setting $K_r = K^{(1)}_r K^{(2)}_r$, we get the result.
\end{proof}

\end{document}
